% !LW recipe = latexmk (xelatex)
\documentclass[11pt,spanish]{article}
\usepackage{portfolio-cv}
\hypersetup{pdftitle={David Barragán | Currículum},pdflang={es}}

\begin{document}
\cvHeader{Ingeniero de Software \cvSeparator Desarrollador Full Stack}{Portafolio}

\textit{Desarrollador Full Stack enfocado en servicios backend con Java y Spring Boot. Experiencia en sistemas ERP, microservicios, aplicaciones web y móviles, despliegues en Azure, CI/CD y análisis de datos con Power BI. Aplico buenas prácticas de desarrollo y principios de código limpio para construir sistemas mantenibles y fáciles de probar.}

\cvSection{Experiencia profesional}
\cvEntry{Kronhos ERP \cvSeparator Desarrollador Full Stack}{dic. 2025 -- jun. 2026}
\begin{cvList}
  \item Desarrollé funcionalidades del ERP en frontend y backend. Implementé APIs REST con Java y Spring Boot para conectar las pantallas del ERP con la lógica de negocio y los datos.
  \item Optimicé el rendimiento de la aplicación y apoyé una migración de base de datos en SQL Server.
  \item Trabajé con los servidores y entornos de aplicación del ERP en Azure.
  \item Creé reportes y paneles en Power BI para explorar los datos del ERP y apoyar el análisis de negocio.
\end{cvList}

\cvEntry{EEASA \cvSeparator Desarrollador backend (freelance)}{2025 -- 2026}
\begin{cvList}
  \item Refactoricé la lógica de incidencias de medidores eléctricos en Spring Boot para facilitar su mantenimiento.
  \item Reestructuré las entidades de incidencias en Oracle y ajusté las reglas de validación de sus registros.
  \item Implementé pantallas de gestión de incidencias con Angular y React.
\end{cvList}

\cvSection{Proyectos seleccionados}
\cvEntry{\cvFuelIQ{} \cvSeparator Control de combustible, viajes y rutas en Azure AKS}{}
\begin{cvList}
  \item Microservicios NestJS con PostgreSQL por servicio y despliegue con Docker y GitHub Actions.
\end{cvList}

\cvEntry{\cvHospital{Suite de Gestión Hospitalaria}}{}
\begin{cvList}
  \item Microservicios con Spring Boot, Eureka y Clean Architecture, PostgreSQL distribuido con Citus e interfaz web en React y TypeScript.
\end{cvList}

\cvEntry{\cvForis{} \cvSeparator Foro comunitario y chat en tiempo real}{}
\begin{cvList}
  \item Monorepo con React, NestJS y PostgreSQL. Publicaciones por tema y análisis estático con SonarQube.
\end{cvList}

\cvSection{Competencias}
\cvDetail{Backend}{Java, Spring Boot, NestJS, APIs REST, microservicios}
\cvDetail{Frontend}{React, TypeScript, Angular, React Native}
\cvDetail{Arquitectura y calidad de código}{Clean Architecture, DDD, JUnit, Playwright, SonarQube}
\cvDetail{Cloud y DevOps}{Azure, Docker, Kubernetes (AKS), Git, GitHub Actions, Jenkins, CI/CD}
\cvDetail{Datos y BI}{SQL Server, PostgreSQL, Oracle, Power BI}
\cvDetail{Habilidades blandas}{Comunicación y liderazgo, documentación técnica y coordinación de tareas en Scrum.}
\cvDetail{Idiomas}{Español nativo \cvSeparator Inglés B2+}

\cvSection{Educación}
\cvEntry{Universidad Técnica de Ambato}{2022 -- Actualidad}
Ingeniería en Software

\cvSection{Cursos y certificaciones}
\begin{cvList}
  \item \cvAWSPractitioner{} \cvSeparator AWS Training and Certification
  \item \cvAWSSimuLearn{} (insignia de formación) \cvSeparator AWS Training and Certification
  \item Curso de Programación de Microcontroladores ESP32 \cvSeparator Universidad Técnica de Ambato, 2025
  \item Desarrollo Web con Python y Django \cvSeparator Universidad Técnica de Ambato
\end{cvList}
\end{document}
